\documentclass[12pt]{article}
\usepackage[utf8]{inputenc}
\usepackage[T1]{fontenc}
\usepackage{hyperref}
\title{Zhire Language Resources for K-12 Curriculum Development}
\author{Compiled by Claude Sonnet 5.5 (claude-sonnet-5-5) for Andr\'as Kornai}
\date{}
\begin{document}
\maketitle

\section{Introduction}
Zhire (ISO 639-3 zhi; Glottocode zhir1238; also called Kenyi, Banyi, Jabji and Kanyi) is a poorly documented Plateau language of Kaduna State, Nigeria. It is spoken in Kagarko LGA. Its closest relative is Shang, which is described as relexified Zhire. Joshua Project reports about 18,000 speakers, far fewer than the figure in the task description. The community is mostly Christian and farms rice, maize and fonio.

\section{Primary Educational Resources}
No primers, textbooks or school programmes in Zhire were found. Zhire is not known to be taught in schools.

\section{Linguistic and Reference Materials}
\begin{itemize}
\item Wikipedia article on the Zhire language: \url{https://en.wikipedia.org/wiki/Zhire_language}
\item Glottolog record for Decker, Dabet, Muniru, Abraham and Jonah (2021), a sociolinguistic profile of Zhire in the Journal of Language Survey Reports: \url{https://glottolog.org/resource/reference/id/585645}
\end{itemize}

\section{Religious and Community Resources}
\begin{itemize}
\item Joshua Project profile of Zhire in Nigeria, with population and religion data: \url{https://joshuaproject.net/people_groups/20774}
\end{itemize}
No Zhire scripture translation or community organisation was found.

\section{Digital and Audio Resources}
\begin{itemize}
\item Global Recordings Network page for Zhire, with audio and slideshow resources: \url{https://globalrecordings.net/en/language/zhi}
\end{itemize}

\section{Curriculum Development Context}
Zhire is used as a first language by the whole community, so it is stable. Written materials are almost absent. A curriculum would first need an orthography, a lexicon and basic literacy materials developed with the community. The 2021 sociolinguistic profile would be a natural starting point.

\end{document}